\subsection{\quantmsdiann{} pipeline architecture}
\label{sec:architecture}

Built on the nf-core framework~\cite{ewels2020nfcore}, \quantmsdiann{} accepts a single SDRF file as primary input, which encodes per-sample metadata (cell line or tissue origin, factor values, instrument, post-translational modifications) and per-run file references. We extended the SDRF with DIA-relevant columns (MS1 and MS2 mass windows, and plexDIA channel definitions and labels) and added an adapter to sdrf-pipelines~\cite{dai2021sdrf} (\texttt{convert-diann}) that translates each run's declared metadata into a per-run \texttt{DIA-NN} command-line configuration. Every run is therefore parametrised from its own metadata with no manual configuration or parameter-harmonisation step, removing a common source of analysis-script divergence between groups: different MS1/MS2 window schemes can be mixed within a cohort, or a plexDIA dataset run with its channel definitions, without any workflow change beyond the SDRF.
The workflow comprises four mandatory and three optional modules organised around the \texttt{DIA-NN} analysis core (Fig.~\ref{fig:architecture}a), alternating between per-file parallel stages (red; raw-file preparation, preliminary calibration, and individual search) and collective stages (green; in-silico library generation, empirical library assembly, final quantification, MSstats conversion, and \texttt{pmultiqc}~\cite{larrea2025pmultiqc} quality-control reporting) operating on the full cohort.

\quantmsdiann{} follows \texttt{DIA-NN}'s distributed multi-pass flow, alternating parallel per-file work with serial cohort-wide steps (Fig.~\ref{fig:architecture}a). Each raw file is first prepared \emph{in parallel} (vendor-format conversion, decompression, or indexing) where needed; conversion to mzML is mandatory for early \texttt{DIA-NN} versions and for some vendor formats such as SCIEX \texttt{.wiff}. A spectral library is then predicted in-silico from the FASTA once (\emph{serial}; skipped when a library is supplied), and every file is calibrated against it in a \emph{parallel} first pass. The first-pass results are merged into a single empirical library (\emph{serial}); each file is searched again against that library \emph{in parallel}; and a final consolidation pass merges all per-file results into the cohort quantification, applying match-between-runs, normalisation, and protein inference (\emph{serial}). MSstats conversion, \texttt{pmultiqc} reporting, and \qpx{} export then run once at the end (\emph{serial}). The two per-file search passes carry essentially all of the compute and scale across cluster nodes, driving the throughput reported below (Section~\ref{sec:throughput}); the serial steps set a largely fixed wall-clock floor.
Outputs include the standard \texttt{DIA-NN} report tables: \texttt{report.tsv} (or \texttt{report.parquet}), \texttt{report.pr\_matrix.tsv}, and \texttt{report.pg\_matrix.tsv}, alongside MSstats-formatted input and a \texttt{pmultiqc}~\cite{larrea2025pmultiqc} quality-control report.
A harmonised \qpx{} archive bundles identifications, metadata, and provenance into a Parquet-based artefact, queried directly with DuckDB and, for the quantification layers, exported to MuData (\texttt{.h5mu}) for scverse-compatible~\cite{virshup2023scverse} downstream analysis.

\subsection{quantmsdiann throughput and scaling}
\label{sec:throughput}

To characterise how \quantmsdiann{} scales on an HPC cluster, we simulated the reanalysis of an entire single-cell experiment of \textasciitilde 2{,}300 raw files (PXD071075) across cluster sizes from 10 to 300 nodes (set through the Nextflow \texttt{executor.queueSize}), and measured the end-to-end wall-clock time (Fig.~\ref{fig:architecture}b). Because the per-file work is parallelised across nodes, a single SDRF-driven invocation completed the whole cohort within hours: wall-clock fell monotonically from 37.4~h at 10 nodes to 2.2~h at 300 nodes (excluding the one-time in-silico library-prediction step, which is independent of cluster size). The curve has a practical knee near 200 nodes (2.4~h), beyond which the serial cohort-level steps dominate and additional nodes give diminishing returns (a 30$\times$ increase in cluster width yields only a 17.4$\times$ speed-up), so the cluster should be sized to the desired turn-around time rather than to maximum throughput.

We then compared the time-to-finish across every dataset in this study (Fig.~\ref{fig:architecture}c), reporting each run net of the one-time in-silico library-prediction step (a cohort-independent cost that can be predicted once and reused). For proteome-scale cohorts, wall-clock stayed within a few hours irrespective of cohort size, because the per-file passes absorb additional files in parallel: the largest cohort (ProCan, PXD030304; 5{,}798 TripleTOF~6600 runs) finished in 6.2~h, in the same few-hour band as cohorts orders of magnitude smaller, while the 6-run ProteoBench benchmarks completed in under 1.5~h. Time-to-finish is therefore governed by the per-file search depth and the remaining cohort-level consolidation rather than by the number of files: the single longest run is a 10-file phosphoproteomics cohort (PXD049692, 8.0~h after excluding library prediction), where the deep variable-modification search dominates the residual wall-clock. Peak memory during the cohort-level consolidation stayed within the 64--128~GB available on commodity HPC nodes, so no high-memory specialist hardware is required. Per-step wall-clock and resource envelopes are reported in Additional file~1: Figs.~S1 and~S2.

\subsection{ProteoBench benchmark across \texttt{DIA-NN} versions and instruments}
\label{sec:proteobench}

To test whether \quantmsdiann{} preserves \texttt{DIA-NN}'s analytical performance under SDRF-driven orchestration, we processed every public ProteoBench DIA module through the workflow and compared the resulting identification counts against the ProteoBench public-submission snapshot (29 May 2026)~\cite{devreese2025proteobench}.
Because \quantmsdiann{} runs \texttt{DIA-NN} library-free (predicting the library in-silico from the FASTA, with no externally supplied experimental library), we restricted the community set for a like-for-like comparison to ProteoBench \texttt{DIA-NN} submissions that used the same predicted-from-FASTA library (\texttt{predictors\_library}~$=$~DIANN); submissions that loaded a pre-existing experimental or user-curated spectral library search a different candidate space and are excluded from the headline comparison.
The benchmark panel covers four instrument families and acquisition modes: Orbitrap Astral 2~Th DIA (\emph{Module~7}), single-cell DIA on Orbitrap Astral (\emph{Module~9}, PXD049412), timsTOF diaPASEF (\emph{Module~5}, PXD062685), and SCIEX ZenoTOF~7600 (\emph{Module~10}, PXD070049).
For each module we ran the workflow against three \texttt{DIA-NN} configurations --- the oldest supported release (1.8.1) and the current release in its academic (2.5.1) and enterprise (2.5.1-enterprise) builds --- in a single invocation, all parametrised from the same SDRF (Fig.~\ref{fig:proteobench}). Identification counts are taken from the per-precursor \texttt{DIA-NN} report rather than the quantities matrices. Each release is run at \texttt{DIA-NN}'s recommended run-specific precursor q-value (\texttt{--qvalue}: 1\% for 1.8.1, 5\% for the 2.5.1 builds); on top of this we apply, uniformly across configurations, a 1\% global precursor q-value, and count unique precursors identified in at least one run. Protein groups are summarised --- like-for-like across versions, each at a 1\% run-specific protein-group q-value --- as the per-run average and the count quantified in \emph{every} run (``complete profiles''); the per-run depth metric tracks a release's gains more directly than the cross-run union, which saturates as rare identifications accumulate across replicates. The recommended run-specific precursor cut-off mainly affects the reproducibility-filtered ($\geq$3-replicate) precursor counts; the single-replicate precursor totals and the per-run protein metrics apply the same thresholds to every version and are directly comparable across releases.

Precursor counts (identified in at least one run) rose from 1.8.1 to 2.5.1 on every module, by 12\% on Module~7 (116{,}771~$\to$~130{,}246), 11\% on Module~9 (41{,}756~$\to$~46{,}243), 21\% on Module~5 (97{,}149~$\to$~117{,}396), and 16\% on Module~10 (88{,}179~$\to$~102{,}457); the enterprise build adds a further 2--3\% on each module.
Protein-group depth rose in step. Averaged per run, protein groups at 1\% FDR increased from 1.8.1 to 2.5.1 by 8--11\% (Module~7 10{,}091~$\to$~11{,}167; Module~9 7{,}157~$\to$~7{,}895; Module~5 10{,}199~$\to$~11{,}231; Module~10 8{,}788~$\to$~9{,}470), with the enterprise build a further 1--3\% higher and highest on every module. The improvement is larger for proteins quantified in \emph{every} run (complete profiles): the enterprise build gains 12--20\% over 1.8.1 (Module~7 8{,}799~$\to$~10{,}403, $+$18\%; Module~9 6{,}124~$\to$~7{,}349, $+$20\%; complete-profile panel in Additional file~1).
Reproducibility-filtered precursor counts ($\geq$3 replicates) follow the same ordering across the three configurations on all four modules.
This monotonic version-progress trend is consistent with what the broader ProteoBench community reports across \texttt{DIA-NN} releases~\cite{devreese2025proteobench}, indicating that the SDRF wrapper introduces no version-coupling artefacts.
Quantification accuracy, by contrast, is essentially invariant across \texttt{DIA-NN} versions.
On every module the measured per-species log$_2$ fold-changes recover the expected HYE ratios (human at unity, yeast and \emph{E.\ coli} near their design ratios, with the mild attenuation at the extreme \emph{E.\ coli} ratio that is characteristic of DIA quantification), and the three \texttt{DIA-NN} configurations overlie one another at every species (Fig.~\ref{fig:proteobench}b).
The version-to-version shift in median quantification error~$|\epsilon|$ ($\sim$0.01 log$_2$, $\approx$0.5--1\% in linear fold-change) is far smaller than the per-precursor spread, so the choice of \texttt{DIA-NN} release is best driven by the identification gains in panel~(a) rather than by accuracy.
On the two modules for which the snapshot contained predicted-library \texttt{DIA-NN} submissions, Orbitrap Astral (Module~7, $n=15$) and timsTOF diaPASEF (Module~5, $n=8$), \quantmsdiann{}'s three configurations fall within the community median-$|\epsilon|$ distribution, at or below its median (Fig.~\ref{fig:proteobench}c); the single-cell and ZenoTOF modules carried no predicted-library community submission. Per-version, per-module protein-group and precursor breakdowns are provided in Additional file~1: Fig.~S3.

\subsection{Independent reanalyses of public DIA deposits recover comparable or deeper coverage than the original quantifications}
\label{sec:reanalyses}

To evaluate \quantmsdiann{} on production-scale public deposits, we reanalysed independent DIA cohorts, spanning bulk cell lines, single cells, and spatial proteomics, that had been processed and published with different DIA engines (Fig.~\ref{fig:cell-line-reanalysis}).
For each cohort the same raw files were retrieved from PRIDE Archive~\cite{deutsch2026proteomexchange} (or MassIVE, for the single-cell plexDIA cohort), parsed through an SDRF generated from the original sample annotations, and run through \quantmsdiann{} with current \texttt{DIA-NN} releases.

On the \textbf{NCI-60} panel (PXD003539, Guo \textit{et al.}~2019~\cite{guo2019nci60}; originally analysed with OpenSWATH against the Guo assay library), \quantmsdiann{} recovered 95{,}633 peptides at 1\% FDR against 40{,}592 in the original submission ($+135$\%), with comparable protein-group counts (6{,}747 versus 6{,}556; Fig.~\ref{fig:cell-line-reanalysis}a).
On the \textbf{single-cell oocyte plexDIA} cohort (MSV000093870, Galatidou \textit{et al.}~2024~\cite{galatidou2024oocyte}; mTRAQ 3-channel, Q~Exactive), the first plexDIA dataset processed through \quantmsdiann{}, the workflow quantified 2{,}904 protein groups from the same raw data, versus 2{,}122 in the original analysis ($+37$\%). At the protein-group level it recovered 94\% of the originally reported groups (1{,}999 of 2{,}122) and detected 905 further groups, at comparable per-cell median depth (1{,}736 versus 1{,}784, across 107 and 97 cells respectively; Fig.~\ref{fig:cell-line-reanalysis}b). A QC-matched per-oocyte comparison (Additional file~1) confirms equivalent single-cell depth (Pearson $r=0.98$).
On the \textbf{ProCan}-DepMapSanger 949-cell-line pan-cancer panel (PXD030304, Gon\c{c}alves \textit{et al.}~2022~\cite{goncalves2022procan}; originally analysed with \texttt{DIA-NN}~1.7.x), \quantmsdiann{} recovered 8{,}921 protein groups at 1\% FDR against 8{,}498 originally ($+5$\%) and 7{,}990 protein groups requiring $\geq 2$ unique peptides versus 6{,}692 ($+19$\%; Fig.~\ref{fig:cell-line-reanalysis}c).
Both original counts are reported in the ProCan paper itself, and the \quantmsdiann{} numbers use the matching ProCan-paper filters under our conservative contaminant/decoy policy (Experimental Procedures), so each side of the comparison is like-for-like.
A fourth, bulk cohort already analysed with a current DIA engine, the Sun \textit{et al.}~2023 PCT-SWATH breast-cancer panel (PXD004701, 76 lines~\cite{sun2023breast}), yielded near-parity (6{,}146 versus 6{,}091 protein groups, $+0.9$\%; Additional file~1), as expected when the original analysis already used a modern engine.
On a \textbf{spatial} Deep Visual Proteomics cohort of MYCN-amplified neuroblastoma (PXD064049; laser-microdissected regions, diaPASEF, originally analysed with \texttt{DIA-NN}~1.8.1), \quantmsdiann{} (\texttt{DIA-NN}~2.5.0) matched the original at the precursor level (16{,}443 versus 16{,}518) and quantified somewhat fewer protein groups (2{,}425 versus 2{,}882; Fig.~\ref{fig:cell-line-reanalysis}d); this protein-group difference reflects \quantmsdiann{}'s entrapment-augmented search, which imposes tighter error control than the original plain-FASTA analysis (only 0.7\% of \quantmsdiann{}'s accepted protein groups map to entrapment sequences), rather than reduced sensitivity.
Finally, on a \textbf{phosphoproteomics} cohort (PXD049692; NK-cell Fe-NTA enrichment, diaPASEF, originally analysed with Spectronaut directDIA), \quantmsdiann{} recovered a comparable number of phosphopeptide backbones at 1\% FDR (distinct stripped sequences: 4{,}234 versus 4{,}254; Fig.~\ref{fig:cell-line-reanalysis}e). At the site-resolved phosphopeptidoform level the two analyses diverge (4{,}592 versus 7{,}993), with Spectronaut reporting more positional isomers; we therefore frame this as comparable phosphopeptide-backbone coverage rather than site-level parity.
Per-condition completeness, peptide-per-protein distributions, gene-versus-accession overlap and per-cell-line missingness for the cell-line and single-cell cohorts are provided as Additional file~1: Figs.~S4--S12.

\subsection{Pan-cohort of five public DIA cell-line reanalyses}
\label{sec:pancohort}

To illustrate the value of running the same SDRF-driven workflow across many independent deposits, we combined the three reanalysed cohorts above with two additional cell-line DIA panels (PXD041421, the MultiPro batch-effect testbed, Wang \textit{et al.}~2023~\cite{wang2023multipro}; PXD017199, Tognetti \textit{et al.}~2021~\cite{tognetti2021breast}, 67 breast-cancer lines) into a single \quantmsdiann{} invocation and assembled a pan-cohort view of per-cohort yield and per-tissue coverage (Fig.~\ref{fig:pancohort}).
Per-cohort headline counts are summarised in Fig.~\ref{fig:pancohort}a.
Across the five cohorts the reanalysis spans 12{,}229 unique UniProt accessions, of which 5{,}144 (42.1\%) are detected in all five (a pan-cohort core proteome) and a further 1{,}334 (10.9\%) in exactly four; these overlap counts are tabulated in the combined per-cohort counts provided with the analysis code.

Tissue-level coverage across the combined panel spans 27 Cellosaurus-derived tissue categories (Fig.~\ref{fig:pancohort}b,c): lung, breast, haematopoietic-and-lymphoid, skin, and central nervous system together account for 59\% of all cell-line entries in the pan-cohort.
Because every count comes from a single SDRF-parameterised \quantmsdiann{} invocation, the resulting protein matrix is queryable from the published \qpx{} archive and can be filtered, joined, or re-grouped by Cellosaurus, NCIt, or Disease Ontology terms without manual harmonisation.
